\section{ضرورت و پیشینه}
\subsection{رمزنگاری پساکوانتومی}
رایانهٔ کوانتومی به‌اندازهٔ کافی بزرگ می‌تواند الگوریتم‌های کلید عمومی رایج مانند \lr{RSA} و منحنی‌های بیضوی را بشکند\cite{shor}. نگرانی عملی «ذخیره اکنون، رمزگشایی بعداً» است: دادهٔ حساسی که امروز رهگیری می‌شود، ممکن است در آینده شکسته شود. به همین دلیل \lr{NIST} در اوت $2024$ نخستین استانداردهای پساکوانتومی را منتشر کرد: \lr{FIPS 203} (\lr{ML-KEM})، \lr{FIPS 204} (\lr{ML-DSA}) و \lr{FIPS 205} (\lr{SLH-DSA})\cite{nist}. مهاجرت سازمان‌ها به این الگوریتم‌ها نیازمند نیروی متخصص برای ارزیابی عملکرد، یکپارچه‌سازی و امنیت پیاده‌سازی است.

\subsection{توزیع کلید کوانتومی}
\lr{QKD} کلید مشترک را با تکیه بر قوانین فیزیک کوانتومی توزیع می‌کند. دو خانوادهٔ اصلی وجود دارد: متغیر گسسته (آشکارساز تک‌فوتون) و متغیر پیوسته (آشکارساز همدوس)\cite{pirandola,weedbrook}. نوع پیوسته با قطعات مخابراتی رایج کار می‌کند و به آشکارساز تک‌فوتون نیاز ندارد\cite{grosshans,diamanti}؛ از این رو برای ستاپ آزمایشی، هم ارزان‌تر و هم از نظر تأمین ساده‌تر است. آزمایش‌های اخیر پیاده‌سازی با امنیت ترکیب‌پذیر و همزیستی با ترافیک کلاسیک را نشان داده‌اند\cite{composable,coexist}.

\subsection{سنجش و رادار کوانتومی}
ایدهٔ روشن‌سازی کوانتومی آن است که از جفت حالت درهم‌تنیده استفاده شود تا حضور هدف در محیط پرنویز آشکار شود\cite{lloyd}. تحلیل نظری برای حالت‌های گاوسی، مزیتی تا $6$ دسی‌بل در نمای خطای آشکارسازی گزارش می‌کند\cite{tan}، و آزمایش‌های نوری و مایکروویو در مقیاس آزمایشگاهی انجام شده‌اند\cite{zhang,barzanjeh}. بهره‌گیری کامل از این مزیت به گیرنده‌های بهینه نیاز دارد که خود موضوعی پژوهشی و جذاب است. این پروپوزال این حوزه را به‌عنوان پژوهش پیشرو و مرحله‌ای تعریف می‌کند.

